\documentclass[12pt]{article}
\usepackage{geometry}
\usepackage{titling}
\usepackage{fontspec}
\usepackage{amsmath,amssymb,amsfonts}
\usepackage{unicode-math}
\usepackage{longtable}
\usepackage{booktabs}
\usepackage{xcolor}
\usepackage{nicematrix}
\usepackage{graphicx}
\usepackage[ruled,linesnumbered,noend,vlined]{algorithm2e}
\usepackage[colorlinks=true, linkcolor=blue!50!black, citecolor=blue!50!black, urlcolor=blue!50!black]{hyperref}

\geometry{a4paper, top=2cm, bottom=2cm, left=2.5cm, right=2.5cm}
\setlength{\droptitle}{-2cm}

\SetKwInOut{Input}{Input}
\SetKwInOut{Output}{Output}
\SetKwInOut{Hyperparameters}{Hyperparameters}
\SetKwInOut{Parameters}{Parameters}
\SetKwInput{KwParams}{Parameters}
\SetKw{KwReturn}{return}
\SetKwComment{Comment}{\big[\!\big[}{\big]\!\big]}
\SetKwFor{For}{for}{do}{end}
\SetKwFor{ForPlain}{for}{}{}
\SetKwInput{KwIn}{Input}
\SetKwInput{KwOut}{Output}
\SetKwInput{KwParams}{Parameters}
\SetKw{KwReturn}{return}
\DontPrintSemicolon
\SetNoFillComment

%\setmainfont{Libertinus Serif}
%\setsansfont{Libertinus Sans}[scale=0.5]
%\setmonofont{DejaVu Sans Mono}[scale=0.5]
\setmathfont{STIXTwoMath-Regular.otf}
%\setmathfont{STIX Two Math}

%
%\newcommand{\scaletop}[1]{\vcenter{\hbox{\scalebox{0.75}{$#1\Trans$}}}}
\newcommand{\scaletop}[1]{\raisebox{1.5pt}{\scalebox{0.75}{$#1\top$}}}
\newcommand{\Trans}{{\mathpalette\scaletop\relax}}
%%
\newcommand{\R}{\mathbb{R}}
\newcommand{\Matrix}[1]{\textbf{#1}}
\newcommand{\Vector}[1]{\textbf{#1}}
\newcommand{\MMat}[1]{\mathbf{#1}}
\newcommand{\MVec}[1]{\mathbf{#1}}
\newcommand{\Norm}[1]{\lVert {\mathbf{#1}} \rVert}
\newcommand{\NormBeta}[1]{\lVert \beta \mathbf{#1} \rVert}
\newcommand{\NormConst}[1]{\lVert #1 \rVert}
\newcommand{\abs}[1]{\lvert {#1} \rvert}
\newcommand{\Avg}[1]{\mathbf{avg}(\mathbf{#1})}
\newcommand{\Std}[1]{\mathbf{std}(\mathbf{#1})}
\newcommand{\Draft}[1]{\big[\!\big[{#1}{\big]\!\big]}}

\let\oldnl\nl % used in the next line.
\newcommand{\nonl}{\renewcommand{\nl}{\let\nl\oldnl}}

\title{Formal Algorithms for Transformers\\[4pt]
    {\large{Multi-Head Attention, Decoder-Transformer and Decoder-Inference\\using\vspace*{-7pt}\\GPT-2 Small Hyperparameters}}
     \\{\small\Draft{Draft 3}}}
\author{Devi Prasad M}
\date{September, 2026}

\begin{document}
\maketitle

\section*{The Formal Algorithms for GPT-2 Small Model}
The technical report Formal Algorithms for Transformers
by Phuong and Hutter of DeepMind (2022, arXiv:2207.09238)
provides a formal yet readable breakdown of the original
Transformer architecture. By explicitly defining matrix
dimensions and tensor shapes in pseudocode, it offers a
clear and succinct description of transformer architectures.
In this report, we add a few extra annotations to their
excellent pseudocode, making it even more accessible to a
broader audience. Thiir work provides the baseline for
the decoder-only models analyzed in this report.

Phuong and Hutter's pseudocode explicitly shows how a
Transformer processes different mathematical objects when
predicting the next token. Despite its clarity, it can
still be difficult to clearly visualise the concrete
shape and structure of the computation. We found that
plugging in specific values for all hyperparameters
makes these operations much easier to track.
Additionally, using an input sequence of a fixed
length helps reveal the inner workings and dynamics
of the attention mechanism.


In this work, we focus on the following algorithms
from the original report, instantiating them with
the hyperparameters of OpenAI's GPT-2 Small model.
We retain the original algorithm names and numbering
scheme for easy cross-referencing:
\begin{itemize}
    \item Algorithm 1: Token embedding
    \item Algorithm 3: Basic single-query attention
    \item Algorithm 4: Attention
    \item Algorithm 5: MHAttention
    \item Algorithm 6: layer\_norm
    \item Algorithm 7: Unembedding
    \item Algorithm 10: DTransformer
    \item Algorithm 14: DInference
\end{itemize}

The pseudocode here follows exactly the same syntax as
the original. We introduce only one change to the syntax
of passing hyperparameter arguments to the algorithms.
We will indicate this change when we first encounter
the affected code.

\section*{Hyperparameters of GPT2-small}
\begin{center}
\begin{tabular}{r@{}llll@{}}
\toprule
\hspace*{2.5em}&$\ell_{\text{max}}$       & $1024$ \qquad\qquad & $d_{\text{e}}$ (model width)            & $768$\\
&$L$ (layers)       & $12$                & $d_\text{mlp}$ (MLP width)          & $3072$ \\
&$H$ (heads)        & $12$                & $d_\text{attn}=d_\text{mid}$ (per head)  & $64$ \\
&$N_{\text{V}}$ (vocabulary) & $50{,}257$          & $d_\text{in}$ (per head)            & $= d_{\text{e}} = 768  \hspace*{2.5em}$ \\
&total parameters ($W_u$)& $\approx124.4$M& $d_\text{out}$ (per head)           & $64$\\
\bottomrule
\end{tabular}
\end{center}
$d_\text{attn}=d_\text{mid}=d_{\text{e}}/H=768/12=64$: each of the 12 heads works in a 64-dimensional
query, key, and value space; $H\cdot d_\text{mid}=768=d_{\text{e}}$, so the concatenated heads equal the
model width before $W_o$ is applied.

\section*{Notation Simplification}
To simplify notation, we denote by $\symbfcal{W}$ the entire set of parameters (query, key,
value and output linear projections) required by a multi-head attention layer:
\begin{equation}
\symbfcal{W} := \left(
\begin{array}{l}
    \symbf{W_q^h} \in \mathbb{R}^{d_\text{attn}\times d_\text{x}},\ \symbf{b_q^h} \in \mathbb{R}^{d_\text{attn}},\ h\in[H] \\[4pt]
    \symbf{W_k^h} \in \mathbb{R}^{d_\text{attn}\times d_z},\ \symbf{b_k^h} \in \mathbb{R}^{d_\text{attn}},\ h\in[H] \\[4pt]
    \symbf{W_v^h} \in \mathbb{R}^{d_\text{mid}\times d_z},\ \symbf{b_v^h} \in \mathbb{R}^{d_\text{mid}},\ h\in[H] \\[4pt]
    \symbf{W_o} \in \mathbb{R}^{d_\text{out}\times Hd_\text{mid}},\ \symbf{b_o} \in \mathbb{R}^{d_\text{out}}
\end{array}
\right)
\label{eq:mha-params}
\end{equation}

For GPT-2 small ($H=12$, $d_\text{attn}=d_\text{mid}=64$, $d_\text{x}=d_z=d_\text{out}=768$), the parameter set
$\symbfcal{W}$ of a single multi-head attention layer becomes:
\begin{equation}
\symbfcal{W} := \left(
\begin{array}{l}
    \symbf{W_q^h} \in \mathbb{R}^{64\times768},\ \symbf{b_q^h} \in \mathbb{R}^{64},\ h\in[12] \\[4pt]
    \symbf{W_k^h} \in \mathbb{R}^{64\times768},\ \symbf{b_k^h} \in \mathbb{R}^{64},\ h\in[12] \\[4pt]
    \symbf{W_v^h} \in \mathbb{R}^{64\times768},\ \symbf{b_v^h} \in \mathbb{R}^{64},\ h\in[12] \\[4pt]
    \symbf{W_o} \in \mathbb{R}^{768\times(12\cdot64)}=\mathbb{R}^{768\times768},\ \symbf{b_o} \in \mathbb{R}^{768}
\end{array}
\right)
\label{eq:mha-params-gpt2small}
\end{equation}

\setcounter{algocf}{2} % so it is numbered "Algorithm 3"
\begin{algorithm}
\caption{Basic single-query attention.}
\DontPrintSemicolon
\KwIn{$\symbf{e} \in \mathbb{R}^{d_{\text{in}}}$, vector representation of the current token}
\KwIn{$\symbf{e}_t \in \mathbb{R}^{d_{\text{in}}}$, vector representations of context tokens $t \in [T]$.}
\KwOut{$\tilde{\symbf{v}} \in \mathbb{R}^{d_{\text{out}}}$, vector representation of the token and context combined.}
\KwParams{$\symbf{W}_{\symbf{q}}, \symbf{W}_{\symbf{k}} \in \mathbb{R}^{d_{\text{attn}} \times d_{\text{in}}}$,
  $\symbf{b}_{\symbf{q}}, \symbf{b}_{\symbf{k}} \in \mathbb{R}^{d_{\text{attn}}}$, the query and key linear projections.}
\KwParams{$\symbf{W}_{\symbf{v}} \in \mathbb{R}^{d_{\text{out}} \times d_{\text{in}}}$,
  $\symbf{b}_{\symbf{v}} \in \mathbb{R}^{d_{\text{out}}}$, the value linear projection.}
$\symbf{q} \leftarrow \symbf{W}_{\symbf{q}}\symbf{e} + \symbf{b}_{\symbf{q}}$\;
$\forall t:\ \symbf{k}_t \leftarrow \symbf{W}_{\symbf{k}}\symbf{e}_t + \symbf{b}_{\symbf{k}}$\;
$\forall t:\ \symbf{v}_t \leftarrow \symbf{W}_{\symbf{v}}\symbf{e}_t + \symbf{b}_{\symbf{v}}$\;
$\forall t:\ \alpha_t = \dfrac{\exp\bigl(\symbf{q}^{\Trans}\symbf{k}_t / \sqrt{d_{\text{attn}}}\bigr)}{\sum_u \exp\bigl(\symbf{q}^{\Trans}\symbf{k}_u / \sqrt{d_{\text{attn}}}\bigr)}$\;
\KwReturn $\tilde{\symbf{v}} = \sum_{t=1}^{T} \alpha_t \symbf{v}_t$\;
\end{algorithm}

\newpage

\section*{List of Notation}
\addcontentsline{toc}{section}{List of notation}
%\begin{group}
%\begin{longtable}{@{} l l p{7cm} @{}}
%\begin{longtable}{l l l p{5cm} l}
\begingroup
\setlength{\tabcolsep}{3pt} % Default is 6pt
\small
\begin{longtable}{@{} p{0.15\textwidth} p{0.15\textwidth} p{0.67\textwidth} @{}}
\toprule
\textbf{Symbol} & \textbf{Type} & \textbf{Explanation} \\
\midrule
\endhead
$[N]$  & $:=\{1,\ldots,N\}$ & set of integers $1,2,\ldots,N-1,N$ \\
$i,j$  & $\in\mathbb{N}$ & generic integer indices \\
$V$    & $\cong[N_{\text{V}}]$ & vocabulary \\
$N_{\text{V}}$  & $\in\mathbb{N}$ & vocabulary size \\
$V^*$   & $=\bigcup_{\ell=0}^{\infty}V^\ell$ & set of token sequences; elements include e.g.\ sentences or documents \\
$\ell_{\text{max}}$ & $\in\mathbb{N}$ & maximum sequence length \\
$\ell$ & $\in[\ell_{\text{max}}]$ & length of token sequence \\
$t$ & $\in[\ell]$ & index of token in a sequence \\
$d_{\ldots}$ & $\in\mathbb{N}$ & dimension of various vectors \\
$\symbf{x}$ & $\equiv x[1\!:\!\ell]$ & $\equiv x[1]x[2]\ldots x[\ell]$ $\in V^\ell$  primary token sequence \\
$\symbf{z}$ & $\equiv z[1\!:\!\ell]$ & $\equiv z[1]z[2]\ldots z[\ell]$ $\in V^\ell$ context token sequence \\
$M[i,j]$ & $\in\mathbb{R}$ & entry $M_{ij}$ of matrix $M\in\mathbb{R}^{d\times d'}$ \\
$M[i,:]\equiv M[i]$ & $\in\mathbb{R}^{d'}$ & $i$-th row of matrix $M\in\mathbb{R}^{d\times d'}$ \\
$M[:,j]$ & $\in\mathbb{R}^{d}$ & $j$-th column of matrix $M\in\mathbb{R}^{d\times d'}$ \\
$\symbf{e}$ & $\in\mathbb{R}^{d_{\text{e}}}$ & vector representation / embedding of a token \\
$\symbf{X}$ & $\in\mathbb{R}^{d_{\text{e}}\times\ell_{\text{x}}}$ & encoded primary token sequence \\
$\symbf{Z}$ & $\in\mathbb{R}^{d_{\text{e}}\times\ell_{\text{z}}}$ & encoded context token sequence \\
$\mathrm{Mask}$ & $\in\mathbb{R}^{\ell_{\text{z}}\times\ell_{\text{x}}}$ & masking matrix, it determines the attention context for each token \\
$L,L_{enc},L_{dec}$ & $\in\mathbb{N}$ & number of network (encoder, decoder) layers \\
$l$ & $\in[L]$ & index of network layer \\
$H$ & $\in\mathbb{N}$ & number of attention heads \\
$h$ & $\in[H]$ & index of attention head \\
$N_{data}$ & $\in\mathbb{N}$ & (i.i.d.) sample size \\
$n$ & $\in[N_{data}]$ & index of sample sequence \\
$\eta$ & $\in(0,\infty)$ & learning rate \\
$\tau$ & $\in(0,\infty)$ & temperature; it controls the diversity-plausibility trade-off at inference \\
$\symbf{W_e}$ & $\in\mathbb{R}^{d_{\text{e}}\times N_{\text{V}}}$ & token embedding matrix \\
$\symbf{W_p}$ & $\in\mathbb{R}^{d_{\text{e}}\times\ell_{\text{max}}}$ & positional embedding matrix \\
$\symbf{W_u}$ & $\in\mathbb{R}^{N_{\text{V}}\times d_{\text{e}}}$ & unembedding matrix \\
$\symbf{W_q}$ & $\in\mathbb{R}^{d_\text{attn}\times d_\text{x}}$ & query weight matrix \\
$\symbf{b_q}$ & $\in\mathbb{R}^{d_\text{attn}}$ & query bias \\
$\symbf{W_k}$ & $\in\mathbb{R}^{d_\text{attn}\times d_z}$ & key weight matrix \\
$\symbf{b_k}$ & $\in\mathbb{R}^{d_\text{attn}}$ & key bias \\
$\symbf{W_v}$ & $\in\mathbb{R}^{d_\text{out}\times d_z}$ & value weight matrix \\
$\symbf{b_v}$ & $\in\mathbb{R}^{d_\text{out}}$ & value bias \\
$\symbfcal{W}_{\symbf{qkv}}$ & & collection of above parameters of a single-head attention layer \\
$\symbf{W_o}$ & $\in\mathbb{R}^{d_\text{out}\times Hd_\text{mid}}$ & output weight matrix \\
$\symbf{b_o}$ & $\in\mathbb{R}^{d_\text{out}}$ & output bias \\
$\symbfcal{W}$ & & collection of above parameters of a multi-head attention layer \\
$\symbf{W}_\text{mlp}$ & $\in\mathbb{R}^{d_1\times d_2}$ & weight matrix corresponding to an MLP layer in a Transformer \\
$\symbf{b}_\text{mlp}$ & $\in\mathbb{R}^{d_1}$ & bias corresponding to an MLP layer in a Transformer \\
$\symbf{\gamma}$ & $\in\mathbb{R}^{d_{\text{e}}}$ & layer-norm learnable scale parameter \\
$\symbf{\beta}$ & $\in\mathbb{R}^{d_{\text{e}}}$ & layer-norm learnable offset parameter \\
$\symbf{\theta},\hat{\symbf{\theta}}$ & $\in\mathbb{R}^d$ & collection of all learnable / learned Transformer parameters \\
\bottomrule
\end{longtable}
\endgroup



\newpage


\medskip
\section*{Mathematical Conventions of the Report}
Phuong and Hutter use the column-vector convention throughout:
tokens are \emph{columns} of $X$ and $Z$, which represent the
primary and the context token sequences, respectively.
The score matrix is $S=K^{\Trans} Q \in \R^{\ell_{\text{z}}\times\ell_{\text{x}}}$
(context length $\times$ primary sequence length), with softmax
taken down each column. This is the transpose
of the row-vector, $S=QK^{\Trans}$. Note that the row-softmax
convention is more commonly used elsewhere including
in the course reference books and in our notes.
The two are mathematically identical, just transposed.
We retain Phuong and Hutter's convention below since our
intention is to match their algorithms exactly.

To begin with consider the expression used to compute an
input mask. This appears in many places in the original
article:
\[
    \operatorname{Mask}[t_{\text{z}}, t_{\text{x}}] =
        \begin{cases}
        1, & t_{\text{z}} \leq t_{\text{x}},\\
        0, & t_{\text{z}} > t_{\text{x}}.
        \end{cases}
\]

We may write the same expression in another way to
convey the fact that the first subscript identifies
rows, and the second, picks columns in a mask matrix:
\[
    \operatorname{Mask}[r,j] =
        \begin{cases}
        1, & r\leq j,\\
        0, & r>j.
        \end{cases}
\]


Recalling that we use the column-vector notation,
it is not difficult to see that a $t^{th} column$
represents a vector whose first $t$ indices are filled
with 1s, followed by zeroes.

\[
    \begin{bNiceMatrix}[first-row,first-col]
        \scriptstyle \symbf{r}\backslash
            \symbf{j} & \scriptstyle 1 & \scriptstyle 2 & \scriptstyle 3 & \scriptstyle 4
                      & \scriptstyle 5 & \scriptstyle 6 & \scriptstyle 7 & \scriptstyle 8\\
        \scriptstyle 1 & 1&1&1&1&1&1&1&1\\
        \scriptstyle 2 & 0&1&1&1&1&1&1&1\\
        \scriptstyle 3 & 0&0&1&1&1&1&1&1\\
        \scriptstyle 4 & 0&0&0&1&1&1&1&1\\
        \scriptstyle 5 & 0&0&0&0&1&1&1&1\\
        \scriptstyle 6 & 0&0&0&0&0&1&1&1\\
        \scriptstyle 7 & 0&0&0&0&0&0&1&1\\
        \scriptstyle 8 & 0&0&0&0&0&0&0&1
    \end{bNiceMatrix}
\]

In the formal algorithms, \(t_{\text{x}}\in[\ell_{\text{x}}]\) indexes positions
of the primary sequence \(\symbf{X}\), whose columns are projected
to form queries, and \(t_{\text{z}}\in[\ell_{\text{z}}]\) indexes positions of the
context sequence \(\symbf{Z}\), whose columns are projected to
form keys and values.


The mask has shape
\(\operatorname{Mask}\in\{0,1\}^{\ell_{\text{z}}\times\ell_{\text{x}}}\).
If \(\operatorname{Mask}[t_{\text{z}},t_{\text{x}}]=1\), the query at position
\(t_{\text{x}}\) may attend to the key–value pair at position \(t_{\text{z}}\).
If it is \(0\), the score \(S[t_{\text{z}},t_{\text{x}}]\) is set to \(-\infty\)
before softmax, so that pair receives zero attention weight.


In self-attention \(\symbf{Z}=\symbf{X}\); the causal mask
\(\operatorname{Mask}[t_{\text{z}},t_{\text{x}}]=[\![t_{\text{z}}\leq t_{\text{x}}]\!]\)
allows each query to attend to its own position and all
earlier positions.

\newpage
\section*{Token Embedding and Unembedding.}
The token embedding learns to represent each
vocabulary element as a vector in $R^{d_{\text{e}}}$.

The unembedding learns to convert a vector
representation of a token and its context into a
distribution over the vocabulary elements.
The algorithm describes an independently learned
unembedding matrix, but note that sometimes
the unembedding matrix is instead fixed to be
the transpose of the embedding matrix.

\setcounter{algocf}{0}
\begin{algorithm}[H]
\caption{Token embedding.}
\label{alg:token_embedding}
\DontPrintSemicolon
    \KwIn{$v \in V \cong [N_{\mathrm{V}}]$, \hfill a token ID.}
    \KwOut{$\symbf{e} \in \mathbb{R}^{d_{\mathrm{e}}}$, \hfill vector representation of the token.}
    \vspace*{2pt}
    \KwParams{$\symbf{W}_{\symbf{e}} \in \mathbb{R}^{d_{\mathrm{e}} \times N_{\mathrm{V}}}$, \hfill token embedding matrix.}
    \BlankLine
    \KwReturn $\symbf{e} = \symbf{W}_{\symbf{e}}[:, v]$\;
\end{algorithm}

\setcounter{algocf}{6}
\begin{algorithm}
    \caption{Unembedding.}
    \label{alg:unembedding}
    \DontPrintSemicolon
    \KwIn{$\symbf{e} \in \mathbb{R}^{d_{\mathrm{e}}}$, \hfill a token encoding.}
    \KwOut{$\symbf{p} \in \Delta(V)$, \hfill a probability distribution over the vocabulary.}
    \vspace*{2pt}
    \KwParams{$\symbf{W}_{\symbf{u}} \in \mathbb{R}^{N_{\mathrm{V}} \times d_{\mathrm{e}}}$, \hfill the unembedding matrix.}
    \BlankLine
    \KwReturn $\symbf{p} = \operatorname{softmax}(\symbf{W}_{\symbf{u}}\symbf{e})$\;
\end{algorithm}


\setcounter{algocf}{5}
\begin{algorithm}
    \caption{$\hat{\symbf{e}} \leftarrow \mathtt{layer\_norm}(\symbf{e} \mid \symbf{\gamma}, \symbf{\beta})$}
    \label{alg:layer_norm}
    \DontPrintSemicolon
    \tcc{Normalizes layer activations $\symbf{e}$.}
    \Input{$\symbf{e} \in \mathbb{R}^{d_{\text{e}}}$, \hfill neural network activations.}
    \Output{$\widehat{\symbf{e}} \in \mathbb{R}^{d_{\text{e}}}$, \hfill normalized activations.}\vspace*{4pt}
    \Parameters{$\symbf{\gamma}, \symbf{\beta} \in \mathbb{R}^{d_{\mathrm{e}}}$, \hfill element-wise scale and offset.}\vspace*{4pt}
    $\symbf{m} \leftarrow \sum_{i=1}^{d_{\text{e}}} \symbf{e}[i] / d_{\text{e}}$\;\vspace*{3pt}
    $v \leftarrow \sum_{i=1}^{d_{\text{e}}} (\symbf{e}[i] - \symbf{m})^2 / d_{\mathrm{e}}$\;\vspace*{3pt}
    \Return $\widehat{\symbf{e}} = \dfrac{\symbf{e} - \symbf{m}}{\sqrt{v}} \odot \symbf{\gamma} + \symbf{\beta}$, \hfill where $\odot$ denotes element-wise multiplication.
\end{algorithm}

\newpage
\section*{Causal (unidirectional/masked) self-attention}

Given a sequence of token representations, Algorithm~\ref{alg:attn}
with the causal mask applies attention at every position, treating
that position and all preceding ones as its context. Because future
positions are masked out, the output at position $t$ depends only on
positions $1,\dots,t$, so this auto-regressive version can be used
for online prediction.

\setcounter{algocf}{3}
\begin{algorithm}[H]
    \caption{$\tilde{\symbf{V}}\leftarrow \textsc{Attention}(\symbf{X}, \symbf{Z} \mid \mathcal{W}_{\symbf{qkv}}, \mathrm{Mask})$}
    \label{alg:attn}
        \Input{$\symbf{X}\in\R^{d_\text{x} \times \ell_\text{x}}$, $\symbf{Z}\in\R^{d_\text{z} \times \ell_\text{z}}$,
               vector representations of the primary and context sequences.\\
               GPT-2: $\symbf{X}=\symbf{Z}$, $d_\text{x} = d_\text{z} = d_\text{e} = 768$, and $\ell_\text{x} = \ell_\text{z}$.}
        \Output{$\tilde{\symbf{V}}\in\R^{d_\text{out} \times \ell_\text{x}}$, updated representations of the tokens in
               $\symbf{X}$, folding in information from the tokens in $\symbf{Z}$.\\
               GPT-2: $d_\text{out} = 64$ (called $d_\text{mid}$ in Algorithm~5).}
        \Hyperparameters{$\mathrm{Mask}\in\{0,1\}^{\ell_\text{z} \times \ell_\text{x}}$, here the causal mask
               $\mathrm{Mask}[t_\text{z},t_\text{x}]=[\![t_\text{z}\le t_\text{x}]\!]$ for
               $t_\text{z}\in[\ell_\text{z}]$, $t_\text{x}\in[\ell_\text{x}]$.}
        \Parameters{$\mathcal{W}_{\symbf{qkv}}$ consisting of:\\
            \quad $\symbf{W}_{\symbf{q}}\in\R^{d_\text{attn} \times d_\text{x}},\;
                   \symbf{W}_{\symbf{k}} \in \R^{d_\text{attn} \times d_\text{z}},\;
                   \symbf{W}_{\symbf{v}} \in \R^{d_\text{out} \times d_\text{z}}$,\\
            \quad $\symbf{b}_{\symbf{q}}\in\R^{d_\text{attn}},\; \symbf{b}_{\symbf{k}}\in\R^{d_\text{attn}},\; \symbf{b}_{\symbf{v}}\in\R^{d_\text{out}}$.\\
            GPT-2: $d_\text{attn} = 64$.}
    \BlankLine
    $\symbf{Q}\leftarrow \symbf{W}_{\symbf{q}} \symbf{X}+\symbf{b}_{\symbf{q}}\symbf{1}^{\Trans}_{\ell_\text{x}}$ \Comment*{Query $\in\R^{d_\text{attn}\times\ell_\text{x}}$}
    $\symbf{K}\leftarrow \symbf{W}_{\symbf{k}} \symbf{Z}+\symbf{b}_{\symbf{k}}\symbf{1}^{\Trans}_{\ell_\text{z}}$ \Comment*{Key $\in\R^{d_\text{attn}\times\ell_\text{z}}$}
    $\symbf{V}\leftarrow \symbf{W}_{\symbf{v}} \symbf{Z}+\symbf{b}_{\symbf{v}}\symbf{1}^{\Trans}_{\ell_\text{z}}$ \Comment*{Value $\in\R^{d_\text{out}\times\ell_\text{z}}$}
    $\symbf{S}\leftarrow \symbf{K}^{\Trans} \symbf{Q}/\sqrt{d_\text{attn}}$ \Comment*{Score $\in \R^{\ell_\text{z} \times \ell_\text{x}}$}
    $\forall t_\text{z}, t_\text{x}$: \textbf{if} $\lnot \mathrm{Mask}[t_\text{z}, t_\text{x}]$ \textbf{then} $\symbf{S}[t_\text{z}, t_\text{x}] \leftarrow -\infty$\;
    \Return{$\tilde{\symbf{V}}=\symbf{V}\cdot\mathrm{softmax}(\symbf{S})$ \Comment*{$\tilde{\symbf{V}} \in\R^{d_\text{out} \times \ell_\text{x}}$}}
\end{algorithm}

In GPT-2 small, each layer has $H = 12$ attention heads. Each head
reads the $768$-dimensional layer-normalized representation of every
token in the input sequence and projects it to its own
$64$-dimensional query, key and value vectors, one of each per token.
In self-attention $\ell_\text{x} = \ell_\text{z} = \ell$, so
$\symbf{Q}^{h}, \symbf{K}^{h}, \symbf{V}^{h} \in \R^{64 \times \ell}$.
The heads \textbf{do not split} the input into twelve $64$-dimensional
pieces; each applies its own learned projections
$\symbf{W}_{\symbf{q}}^{h}, \symbf{W}_{\symbf{k}}^{h},
\symbf{W}_{\symbf{v}}^{h} \in \R^{64 \times 768}$
to the whole representation.

Line~2 of Algorithm~5 calls Algorithm~\ref{alg:attn} once per head,
passing that head's parameters; inside Algorithm~\ref{alg:attn} we
drop the head index $h$. Lines~1--3 apply the affine map
$\symbf{x} \mapsto \symbf{W}_{\symbf{q}}\symbf{x} + \symbf{b}_{\symbf{q}}$
to every column of the primary sequence $\symbf{X}$, giving
$d_\text{attn} = 64$-dimensional queries, and the maps
$\symbf{z} \mapsto \symbf{W}_{\symbf{k}}\symbf{z} + \symbf{b}_{\symbf{k}}$ and
$\symbf{z} \mapsto \symbf{W}_{\symbf{v}}\symbf{z} + \symbf{b}_{\symbf{v}}$
to every column of the context sequence $\symbf{Z}$, giving
$d_\text{attn} = 64$-dimensional keys and $d_\text{out} = 64$-dimensional
values. In GPT-2's self-attention, $\symbf{Z} = \symbf{X}$.

Line~4 divides each query--key dot product by
$\sqrt{d_\text{attn}} = 8$, and line~5 sets the scores at masked
positions to $-\infty$. In line~6, a softmax over the context
positions, taken one column of scores at a time, gives the attention
weights of each query, and the product
$\symbf{V}\cdot\mathrm{softmax}(\symbf{S})$ forms, for each query
token, the weighted average of the value vectors under those weights:
\[
  \tilde{\symbf{V}}[:, t_\text{x}] = \sum_{t_\text{z}=1}^{\ell_\text{z}}
  \mathrm{softmax}(\symbf{S})[t_\text{z}, t_\text{x}]\;\symbf{V}[:, t_\text{z}].
\]
Because the softmax weights are non-negative and sum to $1$, each
output is a convex combination of the value vectors, and since
$\exp(-\infty) = 0$, masked positions receive weight exactly $0$.

Each head therefore outputs one $64$-dimensional vector per query
token. Line~3 of Algorithm~5 stacks the twelve head outputs into a
vector of $H\,d_\text{mid} = 768$ entries for each query token, forming
$\symbf{Y} \in \R^{768 \times \ell_\text{x}}$, and line~4 applies the
output projection:
\[
  \tilde{\symbf{V}} = \symbf{W}_{\symbf{o}}\symbf{Y}
    + \symbf{b}_{\symbf{o}}\symbf{1}_{\ell_\text{x}}^{\Trans},
  \qquad
  \symbf{W}_{\symbf{o}} \in \R^{768 \times 768},\;
  \symbf{b}_{\symbf{o}} \in \R^{768}.
\]
Note the name clash: in Algorithm~5, $d_\text{out} = 768$ is the size
of this final output, whereas in Algorithm~\ref{alg:attn} $d_\text{out}
= 64$ is the size of a single head's output, which Algorithm~5 calls
$d_\text{mid}$.

The multi-head output $\tilde{\symbf{V}} \in \R^{768 \times \ell_\text{x}}$
contains the updated token representations. Each column incorporates
information from the positions that token attends to. Under the causal
mask, $\tilde{\symbf{V}}[:, 1{:}t]$ depends only on $\symbf{X}[:, 1{:}t]$.
Every other operation in the network (the embeddings, residual
connections, layer normalization, MLP, and unembedding) acts on each
column separately. Causal self-attention is thus the only step that
mixes information across positions, and each position attends only to
itself and earlier positions. Applying this reasoning layer by layer,
column $t$ of the final output depends only on the tokens $x[1{:}t]$
and can therefore be used to predict $x[t+1]$.

\subsection*{Prediction and the Causal Mask}
Nothing in the prefix $x[1{:}t]$ determines $x[t+1]$. To
``predict'' $x[t+1]$ means to estimate its conditional distribution
$\hat{P}(x[t+1] \mid x[1{:}t])$ from the observed prefix, not to
determine it with certainty. The knowledge behind this estimate lies
in the trained parameters $\hat{\symbf{\theta}}$, which hold what was
learned from the training corpus (for GPT-2, about 8 million web
documents, WebText): which tokens tend to follow which contexts,
across a very large number of examples. The prefix selects the
context while $\hat{\symbf{\theta}}$ supplies what usually comes
next in such contexts.

The causal mask ensures that this estimate never uses the answer, or
any later token. This matters most in training, where the whole
sequence is present in the input and column $t$ could otherwise read
$x[t+1]$ directly. At inference there is nothing later to hide, but
the mask stays.

\newpage
%\text{the query, key and value projections.}
% $K^{h}\leftarrow W^{h}_{\!\!\!k} \, X+b_k\mathbf{1}^{\Trans}_{\ell_\text{z}}$ \Comment*{Key $\in\R^{d_{\text{attn}}\times\ell_{\text{z}}}$}

\setcounter{algocf}{4}
\begin{algorithm}[H]
\caption{$\tilde{V}\leftarrow \textsc{MHAttention}(X, Z, \mathcal{W}, \mathrm{Mask})$}
\label{alg:mha}
    \Input{$X\in\R^{d_\text{x} \times \ell_{\text{x}}}$ \hfill layer-normalized primary sequence.\\
                $Z\in\R^{d_z \times \ell_{\text{z}}}$ \hfill layer-normalized context sequence.\\
                               \hfill $X = Z$, $\ell_{\text{x}} = \ell_{\text{z}}$}
    \Output{$\tilde{V} \in \R^{d_\text{out} \times \ell_{\text{x}}}$ \hfill attention output.}
    \Hyperparameters{$H=12$; $\mathrm{Mask}\in\{0,1\}^{\ell_{\text{z}} \times \ell_{\text{x}}}$}
    \Parameters{$\mathcal{W}$ consisting of:\\
        \hspace*{-2em}\textbf{for} $h\in[12]$: $W_{qkv}^h$ consisting of:\\
            \hspace*{-2em}\quad $\mid W_q^h \in \R^{d_\text{attn} \times d_\text{x}}, \, b^h_q \in R^{d_\text{attn}}$\vspace*{3pt}\\
            \hspace*{-2em}\quad $\mid W_k^h \in \R^{d_\text{attn} \times d_z}, \, b^h_k \in R^{d_\text{attn}}$\vspace*{3pt}\\
            \hspace*{-2em}\quad $\mid W_v^h \in \R^{d_\text{mid}  \times d_z}, \, b^h_v \in \R^d_\text{mid} $\vspace*{3pt}\\
        \hspace*{-2em}$W_o \in \R^{d_\text{out}\times(h \cdot d_\text{mid})}=\R^{d_\text{out} \times d_\text{out}}$, $b_o\in\R^{d_\text{out}}$, \hfill output projection.}
    \BlankLine
    \For{$h=1,2,\dots,12$}{
    $Y^h\leftarrow \textsc{Attention}(X, X, W_{qkv}^h,\mathrm{Mask})$ \Comment*{Algorithm~\ref{alg:attn}; $Y^h\in\R^{64\times\ell}$}
    }
    $Y\leftarrow[Y^1;Y^2;\dots;Y^{12}]$ \Comment*{stack heads: $Y\in\R^{768\times\ell}$}
    \Return{$\tilde{V}=W_oY+b_o\mathbf{1}^{\Trans}$}
\end{algorithm}

\noindent Note that a column $t$ of the output matrix
$\tilde{V} \in \R^{d_\text{out} \times \ell_{\text{x}}}$
is a convex combination of the value vectors $v_{t'}$
over the permitted positions $t'$.

\newpage
\section*{Decoder-Transformer}

\setcounter{algocf}{9}
\begin{algorithm}[H]
\SetAlgoLined
\SetAlgoLongEnd
\caption{$P\leftarrow \textsc{Decoder-Transformer}(x, \symbf\theta)$}
\label{alg:dtransformer}
    \Input{$x\in V^*$, a sequence of token IDs. \hfill{$length(x) \le \ell_{\text{max}} \le 1024$}}
                    \vspace*{4pt}
    \Output{$P\in(0,1)^{{N_{\text{V}}} \times length(x)}$. \hfill{$N_{\text{V}} = 50257$}}
                    \vspace*{4pt}
    \Hyperparameters{$\ell_{\text{max}},\ L \ H, \ d_{\text{e}},\ d_\text{mlp} \in \mathbb{N}$ \\
                    $\ell_{\text{max}} = 1024,\ L=12 \ H=12, \ d_{\text{e}}=768,\ d_\text{mlp} =3072$ }
                    \vspace*{4pt}
    \Parameters{$\symbf\theta$ includes: \vspace*{2pt}\\
        \hspace*{-4em}$W_e\in\R^{d_{\text{e}} \times N_{\text{V}}}$, \hfill token embeddings.\\
        \hspace*{-4em}$W_p\in\R^{d_{\text{e}} \times \ell_{\text{max}}}$ \hfill positional embeddings.\\
        \hspace*{-4em}{\textbf{for}} $l\in[12]$: \hfill{hyperparameters of each layer $\ldots$}\\
        \hspace*{-4em} \quad $\mid \mathcal{W}_l$, \hfill attention parameters of Algorithm~\ref{alg:mha}.\vspace*{6pt}\\
        \hspace*{-4em} \quad $\mid \ \gamma^{(1)}_l,\beta^{(1)}_l,\gamma^{(2)}_l,\beta^{(2)}_l\in\R^{d_{\text{e}}}$, \hfill two sets of layer-norm parameters.\vspace*{6pt}\\
        \hspace*{-4em} \quad $\mid \ W_{\text{mlp1}}^l\in\R^{d_\text{mlp} \times d_{\text{e}}}$, $b_{\text{mlp1}}^l\in\R^{d_\text{mlp}}$, \hfill MLP parameters.\vspace*{6pt}\\
        \hspace*{-4em} \quad $\mid \ W_{\text{mlp2}}^l\in\R^{d_{\text{e}} \times {d_\text{mlp}}}$, $b_{\text{mlp2}}^l\in\R^{d_{\text{e}}}$, \hfill MLP parameters.\vspace*{6pt}\\
        \hspace*{-4em}$\gamma,\beta\in\R^{d_{\text{e}}}$, \hfill{final layer-norm parameters.}\vspace*{6pt}\\
        \hspace*{-4em}$W_u \in \R^{{N_{\text{V}}} \times d_{\text{e}}}$, \hfill the unembedding matrix.}
    \BlankLine
    $\ell \leftarrow \mathrm{length}(x), \ \, \ell_\text{x} \leftarrow \mathrm{length}(x)$\;
    \ForPlain{$i\in[\ell]$} {
        \nonl $t \leftarrow x[i]$ \Comment*{next token id}
        \nonl $e_i \leftarrow W_e[:,t] + W_p[:,i]$ \Comment*{token embedding; $e_i \in \R^{768}$}
    }
    \LinesNumbered
    $X\leftarrow[e_1,e_2,\dots,e_\ell]$ \Comment*{$X\in\R^{768\times\ell}$}
    \BlankLine
    \For{$l=1,2,\dots,12$} {
        \ForPlain{$t\in[\ell]$} {
            \nonl $\tilde{X}[:,t]\leftarrow\mathtt{layer\_norm}(X[:,t]\mid\gamma^{(1)}_l,\beta^{(1)}_l)$\;
        }
        \vspace*{10pt}
        \tcc{residual attention mixes the unnormalized $x$.}
        $X\leftarrow X+\textsc{MHAttention}(\tilde{X}, \tilde{X}, \mathcal{W}_l,\mathrm{Mask})$ \Comment*{Algorithm~\ref{alg:mha}}
        \ForPlain{$t\in[\ell]$} {
            \nonl $\tilde{X}[:,t]\leftarrow\mathtt{layer\_norm}(X[:,t], \gamma^{(2)}_l,\beta^{(2)}_l)$\;
        }
        \LinesNumbered \vspace*{10pt}
        \tcc{residual attention $x$.}
        $X\leftarrow X+W_{\text{mlp2}}^l\,\mathtt{GELU}\!\left(W_{\text{mlp1}}^l\tilde{X}+b_{\text{mlp1}}^l\mathbf{1}^{\Trans}\right)+b_{\text{mlp2}}^l\mathbf{1}^{\Trans}_{\color{red}{\ell_{\text{x}}}}$ \Comment*{$768\to3072\to768$}
    }
    \BlankLine
    \ForPlain{$t\in[\ell]$} {
        \nonl $X[:,t]\leftarrow\mathtt{layer\_norm}(X[:,t], \gamma,\beta)$\;
    }
    \BlankLine\BlankLine
    \Return{$P=\mathrm{softmax}(W_uX)$}
\end{algorithm}

\newpage
\section*{Decoder-Inference}
The hyperparameters $\ell_{\text{gen}} \in \mathbb{N}$ and
$\tau \in (0,\infty)$ belong to DInference,
not to DTransformer. They are not part of the trained
parameters $\hat{\symbf{\theta}}$ and are chosen at
inference time.

$\ell_{\text{gen}}$ is the number of tokens generated, and hence
the number of calls to DTransformer. Since the last call sees
$\ell + \ell_{\text{gen}} - 1$ tokens, we need
$\ell + \ell_{\text{gen}} - 1 \le \ell_{\max}$.

$\tau$ is the sampling temperature. Each new token is drawn from
$\symbf{q} \propto \symbf{p}^{1/\tau}$, where $\symbf{p}$ is the
last column of $\symbf{P}$. With $\tau = 1$ we sample from
$\symbf{p}$ itself. With $\tau < 1$ the distribution is sharpened,
and as $\tau \to 0$ sampling becomes greedy $\arg\max$ decoding
(it concentrates all its mass on the most probable token).
With $\tau > 1$ the distribution is flattened, approaching
the uniform distribution over tokens as $\tau \to \infty$.

For a fixed input sequence, neither \(\ell_{\text{gen}}\)
nor \(\tau\) changes the \(\symbf{P}\) returned by
DTransformer. $\ell_{\text{gen}}$ determines how many
tokens are generated, and $\tau$ how each next token
is selected. Because a selected token becomes part of the
next input sequence, changing \(\tau\) (or simply a
different random draw) can indirectly change
\(\symbf{P}\) at later generation steps.

\setcounter{algocf}{13}
\begin{algorithm}[H]
\caption{$y\leftarrow \textsc{DInference}(x,\widehat\theta)$}
\label{alg:dinference}
    \Input{$\widehat\theta$, Trained parameters \hfill(shapes as in Algorithm~\ref{alg:dtransformer}).}
    \Input{$x\in V^*$, \hfill a prompt, $\mathrm{length}(x)\le 1024$.}
    \Output{$y\in V^*$, \hfill model's continuation of the prompt.}
    \Hyperparameters{$\ell_\text{gen}\in\mathbb{N}$, $\tau\in(0,\infty)$ \hfill inference-time choices.\\
    \hfill{}}
    \BlankLine
    $\ell\leftarrow\mathrm{length}(x)$\;
    \For{$i=1,2,\dots,\ell_\text{gen}$} {
        $P\leftarrow \textsc{DTransformer}(x, \widehat\theta)$ \Comment*{Algorithm~\ref{alg:dtransformer}}
        $\symbf{p}\leftarrow P[:,\ell+i-1]$\;
        sample a token $y_i$ from $\symbf{q} \propto \symbf{p}^{1/\tau}$\;
        $x\leftarrow[x,y_i]$\;
    }
    \vspace*{2pt}
    \Return{$y=x[\ell{+}1:\ell+\ell_\text{gen}]$}
\end{algorithm}

With fixed $\widehat\theta$, Algorithm 10 is a deterministic
function: the same $x$ always gives the same $P$. Every
line in it is a fixed computation with no randomness:
embedding lookups, layer norm, matrix products, softmax and GELU.

The explicit randomness in Algorithm 14 is in line 5,
where \(y\) is sampled from $y$ is sampled
from $q \propto p^(1/\tau)$. Therefore, repeated runs
from the same prompt can produce altogether different continuations.
If we replace sampling with greedy decoding, DInference
becomes deterministic for fixed parameters and
generation length, provided ties are broken by
a fixed rule.

Further, Algorithm 10 has no dropout, and GPT-2's
published inference code (found in src/model.py) does not
use dropout. In frameworks that include dropout,
it is switched off during inference.

On GPUs, numerical results can sometimes vary between
repeated runs because some operations are nondeterministic.
Results can also differ across hardware, kernels, or
precision settings: floating-point addition is not associative,
so a different order of summation can change the last few bits.
This is an implementation-level effect, not part of the
mathematical algorithm. Even with greedy decoding,
a sufficiently small numerical difference can flip the
argmax between two closely tied tokens.

\newpage

\section*{Parameter-count check}
Summing every shape in Algorithm~\ref{alg:dtransformer}'s
parameter list (with Algorithm~\ref{alg:mha}'s attention parameters,
$12\times(3\times(64{\times}768+64))+768{\times}768+768=23{,}62{,}368$
per layer, expanded in):
\begin{align*}
    & \underbrace{768\,{\times}\,50257}_{W_e} \, + \\[2pt]
    & \underbrace{768\,{\times}\,1024}_{W_p} \, + \\[2pt]
    & 12\,\big(\underbrace{23{,}62{,}368}_{\text{attn/layer}} \, +
    \underbrace{3{,}072}_{\text{LN/layer}} \, +
    \underbrace{47{,}22{,}432}_{\text{MLP/layer}}\big) \, + \\[2pt]
    & \underbrace{1{,}536}_{\text{final LN}} \\[2pt]
    & = 12{,}44{,}39{,}808
\end{align*}
with $W_u$ tied to $W_e$ (untying adds another $768{\times}50257\approx38.6$M, giving $\approx163$M).
$124.4$M matches the commonly cited ``124M'' parameter count for GPT2 small model.

\end{document}
